\section{Introduction}\label{sec:intro}
% src: S-PAPER-1
A Polymarket YES share pays \PayoutDollar{} if the event happens.
At a price of \HookPrice, the price is read as a \HookPct{} chance.
Is that chance right?
If we collect every event priced at \HookPct, about \HookPct{} of them should have happened.
The Polymarket-v1 database records every trade on Polymarket's first exchange.
It holds \PaperTrades{} trades and \PaperMarkets{} markets from \PaperPeriod{} \citep{qin2026polymarketv1}.
So we can check each price against what happened.

This paper reports the first stage of a staged study.
% src: S-PAPER-8
It answers three questions.
The first question is Q1: is the crowd right?
% src: S-PAPER-1
We test calibration one price range at a time.
We use binomial and chi-square tests, confidence intervals, and the mean squared residual of the price.
We also ask whether cheap longshot bets are a trap.
The second question is Q2: who is the crowd?
We ask whether a small group of wallets does most of the trading.
We use top shares of wallet volume, the 80/20 rule, and a power-law check.
We also group wallets by k-means, as a side analysis.
The third question is Q3: what does betting look like?
We describe trade sizes, wallet volumes, and prices with histograms and distribution fits.

% src: S-PAPER-3, S-PAPER-5
Our contributions are these.
We fix the evaluation rules before we see any Q1 result.
We report one rule change made afterwards.
We score all \NMarketsKept{} kept markets, not a small panel.
We report an effect size next to every p-value, such as the per-bin gap.
The text uses the word wallet for an address on the exchange.
